% Lineare Gleichungssysteme – bank/lineare-gleichungssysteme/ (zusammenbau v0.4, Heft)
\einheitenkopf[t4]{Lineare Gleichungssysteme}
\setcounter{aufgabe}{24}

% Anlauf (ohne Prüfkennung)
\begin{aufgabe}{Grafisch – Anlauf}
\begin{teile}
\teil $(2|6)$ – Lösung von welcher Gleichung? I: $x + y = 8$, II: $3x - y = 0$. Kreuze an.\\ \kreuz{nur von I}\\ \kreuz{nur von II}\\ \kreuz{von beiden}\\ \kreuz{von keiner}
\teil $x + y = 5$ – Lösungspaare? Ergänze die Tabelle.

\wertetabelle{x}{y}{0,1,2,3}
\teil $2x + y = 6$ – nach $y$ umstellen. y = \leerfeld
\end{teile}
\end{aufgabe}

\begin{aufgabe}{Grafisch – Prüfungsaufgaben}
\begin{teile}
\steil Zelte: Ein Zeltlager hat $14$ Zelte, Zweierzelte und Viererzelte, mit zusammen $44$ Schlafplätzen. $x$ ist die Zahl der Zweierzelte, $y$ die Zahl der Viererzelte. Wie viele Zelte jeder Art sind es? Löse durch Probieren oder grafisch und begründe, dass es keine andere Möglichkeit gibt. \hfill (P10 2016 OS) \\ x = \leerfeld , y = \leerfeld
\steil Tische: Im Festsaal stehen $15$ Tische, Vierertische und Sechsertische, mit zusammen $74$ Plätzen. $x$ ist die Zahl der Vierertische, $y$ die Zahl der Sechsertische. Wie viele Tische jeder Art sind es? Löse durch Probieren oder grafisch und begründe, dass es keine andere Möglichkeit gibt. \hfill (P10 2016 OS) \\ x = \leerfeld , y = \leerfeld
\end{teile}
\end{aufgabe}

% Anlauf (ohne Prüfkennung)
\begin{aufgabe}{Einsetzen – Anlauf}
\begin{teile}
\teil Welche Gleichung ist fast fertig? I: $3x + 7y = 17$, II: $x + 4y = 9$. Kreuze die Gleichung an, die du umstellst und in die andere einsetzt.\\ \kreuz{Gleichung I}\\ \kreuz{Gleichung II}
\teil I: $y = 2x + 1$, II: $3x + y = 16$ – Lösung? Setze I in II ein. ( \leerfeld | \leerfeld )
\teil I: $y = x + 4$, II: $5x - y = 16$ – Lösung? Setze I in II ein. ( \leerfeld | \leerfeld )
\end{teile}
\end{aufgabe}

\begin{aufgabe}{Einsetzen – Prüfungsaufgaben}
\begin{teile}
\steil Fahrradladen: Ein Helm und ein Schloss kosten zusammen $54{,}00\,€$. Ein Verein kauft $7$ Helme und $4$ Schlösser für $316{,}20\,€$. Mit $x$ als Preis eines Helms und $y$ als Preis eines Schlosses in Euro gilt I: $x + y = 54$, II: $7x + 4y = 316{,}20$. Berechne die Preise und gib an, was ein Helm und was ein Schloss kostet. \hfill (P10 2022 OS) \\ Helm \leerfeld[€] , Schloss \leerfeld[€]
\steil Bootsverleih: Ein Kanu und ein Kajak kosten für einen Tag zusammen $48{,}00\,€$. Eine Gruppe mietet $5$ Kanus und $2$ Kajaks für $199{,}20\,€$. Mit $x$ als Tagespreis eines Kanus und $y$ als Tagespreis eines Kajaks in Euro gilt I: $x + y = 48$, II: $5x + 2y = 199{,}20$. Berechne die Preise und gib an, was ein Kanu und was ein Kajak kostet. \hfill (P10 2022 OS) \\ Kanu \leerfeld[€] , Kajak \leerfeld[€]
\end{teile}
\end{aufgabe}

% Anlauf (ohne Prüfkennung)
\begin{aufgabe}{Sachaufgaben – Anlauf}
\begin{teile}
\teil Bäckerei: $3$ Brötchen und $2$ Brezeln kosten $3{,}10\,€$; $1$ Brötchen und $4$ Brezeln kosten $3{,}70\,€$. Was ist unbekannt? Benenne die gesuchten Größen mit $x$ und $y$ und schreib heraus, welche Zahlen rechts vom Gleichheitszeichen und welche vor $x$ und $y$ stehen.
\teil Ein Eis und ein Kakao kosten zusammen $4{,}60\,€$. Drei Eis und vier Kakao kosten $15{,}80\,€$. $x$: Preis eines Eises, $y$: Preis eines Kakaos (in €). Gleichungen? Stelle I und II auf. I: \leerfeld \quad II: \leerfeld
\teil Freibad: Am Morgen wurden $50$ Karten verkauft, Kinderkarten zu $2{,}80\,€$ und Erwachsenenkarten zu $4{,}50\,€$, zusammen für $177{,}40\,€$. $x$: Anzahl der Kinderkarten, $y$: Anzahl der Erwachsenenkarten. Gleichungen? Stelle I und II auf. I: \leerfeld \quad II: \leerfeld
\end{teile}
\end{aufgabe}

\begin{aufgabe}{Sachaufgaben – Prüfungsaufgaben}
\begin{teile}
\teil Café: Ein Croissant und ein Muffin kosten zusammen $4{,}10\,€$. Fünf Croissants und zwei Muffins kosten $13{,}60\,€$. $x$ ist der Preis eines Croissants, $y$ der Preis eines Muffins in Euro. Stelle ein Gleichungssystem auf. \hfill (P10 2024 OS) \\ I: \leerfeld \quad II: \leerfeld
\teil Baumarkt: Ein Hammer und eine Zange kosten zusammen $23{,}50\,€$. Zwei Hämmer und drei Zangen kosten $56{,}50\,€$. $x$ ist der Preis eines Hammers, $y$ der Preis einer Zange in Euro. Stelle ein Gleichungssystem auf. \hfill (P10 2024 OS) \\ I: \leerfeld \quad II: \leerfeld
\steil Obststand: Äpfel kosten $0{,}60\,€$, Mangos $1{,}40\,€$ das Stück. Mia bezahlt $11{,}60\,€$. Sie schreibt I: $0{,}60a + 1{,}40m = 11{,}60$, II: $a + m = 14$ ($a$ Anzahl der Äpfel, $m$ Anzahl der Mangos). Formuliere II als Satz und löse das System. \hfill (P10 2024 OS) \\ a = \leerfeld , m = \leerfeld
\steil Sportverein: Ein Ball kostet $14{,}50\,€$, ein Hütchen $3{,}50\,€$. Der Verein zahlt $104\,€$. Der Trainer schreibt I: $14{,}50b + 3{,}50h = 104$, II: $b + h = 14$ ($b$ Anzahl der Bälle, $h$ Anzahl der Hütchen). Formuliere II als Satz und löse das System. \hfill (P10 2024 OS) \\ b = \leerfeld , h = \leerfeld
\teil Jugendclub: Ein Zelt und ein Schlafsack kosten zusammen $190\,€$. Der Club kauft $4$ Zelte und $6$ Schlafsäcke für $900\,€$. Dazu gehört I: $x + y = 190$, II: $4x + 6y = 900$. Was bedeuten $x$ und $y$? \hfill (P10 2022 OS)
\teil Schulhof: Eine Tischtennisplatte und ein Kicker kosten zusammen $900\,€$. Die Schule kauft $2$ Platten und $1$ Kicker für $1\,300\,€$. Dazu gehört I: $x + y = 900$, II: $2x + y = 1\,300$. Was bedeuten $x$ und $y$? \hfill (P10 2022 OS)
\teil Musikschule: Jonas bucht $2$ Gitarrenstunden und $3$ Klavierstunden für $94{,}50\,€$, Ella $1$ Gitarrenstunde und $2$ Klavierstunden für $57{,}00\,€$. Stelle zwei Gleichungen auf und berechne den Preis einer Gitarrenstunde und einer Klavierstunde. \hfill (P10 2021 OS) \\ Gitarre \leerfeld[€] , Klavier \leerfeld[€]
\teil Fitnessstudio: Herr Kaya zahlt für $3$ Monatskarten und $4$ Kurskarten $177\,€$, Frau Berg für $1$ Monatskarte und $2$ Kurskarten $71\,€$. Stelle zwei Gleichungen auf und berechne den Preis einer Monatskarte und einer Kurskarte. \hfill (P10 2021 OS) \\ Monatskarte \leerfeld[€] , Kurskarte \leerfeld[€]
\end{teile}
\end{aufgabe}
